\section{Discusión y Explotación Industrial}
\label{sec:discusion}

Los resultados permiten pasar de una comparación bioquímica a una decisión
operativa. La pregunta ya no es solamente qué enzima es mejor, sino qué propiedad
genera valor en cada proceso y en qué condición conviene explotarla.

\subsection{Diagnóstico por Aplicación}
\label{subsec:diagnostico_industrial}

\begin{table}[H]
\centering
\tabcaptionlink{anx:tab_diagnostico}{Diagnóstico Industrial de los Cuatro Casos}{tab:diagnostico_industrial}



\renewcommand{\arraystretch}{1.3}

\begin{tabularx}{\textwidth}{
|p{2.4cm}|
p{3.2cm}|
p{3.2cm}|
X|
}
\hline
Aplicación & Cuello de botella observado & Evidencia principal & Acción prioritaria \\ \hline

Alimentos &
Inactivación temprana &
\(t_{50}\): 1.11 \(\rightarrow\) 3.53 h &
Explotar mayor vida operativa antes de buscar más velocidad. \\ \hline

Bagazo &
Inactivación e inhibición por glucosa &
\(t_{50}\): 0.75 \(\rightarrow\) 7.02 h &
Priorizar estabilidad y tolerancia al producto. \\ \hline

Biomasa &
Ventana térmica &
\(T_{\mathrm{opt}}\): 65 \(\rightarrow\) 75 \(^{\circ}\)C &
Aprovechar temperaturas mayores sin exceder el punto productivo. \\ \hline

Glucósidos &
Inhibición y selectividad &
\(I_C=2.56\), \(I_P=15.29\) &
Optimizar producto útil, no velocidad aislada. \\ \hline

\end{tabularx}
\end{table}

La evidencia gráfica que respalda estos diagnósticos se encuentra en el
\anexolink{anx:estabilidad}{Estabilidad Térmica}, el
\anexolink{anx:temperatura}{Respuesta Térmica} y el
\anexolink{anx:selectividad}{Selectividad en Síntesis de Glucósidos}.

\subsection{Cuánto de la Mejora Llega a Producción}
\label{subsec:traduccion_produccion}

El factor de traducción productiva \(F_P\), definido en la Sección
\ref{subsec:indicadores}, permite distinguir entre mejora molecular y mejora de
proceso.

\begin{table}[H]
\centering
\tabcaptionlink{anx:tab_traduccion}{Traducción de la Mejora Molecular a Producción}{tab:factor_traduccion}



\renewcommand{\arraystretch}{1.25}

\begin{tabular}{|l|c|l|}
\hline
Aplicación & \(F_P\) & Lectura \\ \hline
Alimentos & 0.88 & Ligera atenuación por restricciones del proceso. \\ \hline
Bagazo de caña & 0.87 & La inhibición absorbe parte de la mejora. \\ \hline
Biomasa & 0.70 & La estabilidad deja de ser el único límite. \\ \hline
Glucósidos & 1.87 & Selectividad y tolerancia amplifican el resultado. \\ \hline
\end{tabular}
\end{table}

Biomasa y glucósidos muestran el contraste más útil. En biomasa, continuar
invirtiendo únicamente en estabilidad ofrece retornos técnicos decrecientes. En
glucósidos, propiedades adicionales amplifican la mejora de estabilidad y
eficiencia.

\subsection{Puntos de Operación para Capturar Valor}
\label{subsec:monetizacion}

Para esta etapa se comparó el mejor punto encontrado para ENG con el mejor punto
encontrado para WT dentro de la misma malla evaluada: \(35--90\,^{\circ}\mathrm{C}\),
0.5--80 mM de sustrato y ocho horas de operación.

Esta comparación es distinta de \(I_P\): \(I_P\) usa una condición base común;
el Cuadro \ref{tab:puntos_operacion} permite que cada variante opere en su mejor
punto dentro del rango analizado.

\begin{table}[H]
\centering
\tabcaptionlink{anx:tab_puntos}{Puntos de Operación con Mayor Producto Dentro del Rango Evaluado}{tab:puntos_operacion}



\renewcommand{\arraystretch}{1.35}
\setlength{\tabcolsep}{3.5pt}

\begin{tabularx}{\textwidth}{
|p{2.1cm}|
>{\centering\arraybackslash}p{2.0cm}|
>{\centering\arraybackslash}p{1.5cm}|
>{\centering\arraybackslash}p{2.0cm}|
>{\centering\arraybackslash}p{1.5cm}|
>{\centering\arraybackslash}p{1.5cm}|
X|
}
\hline
Aplicación &
ENG: \(T\), sustrato &
Producto ENG &
WT: \(T\), sustrato &
Producto WT &
Ventaja &
Decisión \\ \hline

Alimentos &
60 \(^{\circ}\)C, 40 mM &
16.47 &
45 \(^{\circ}\)C, 40 mM &
5.93 &
2.78\(\times\) &
Usar ENG para extender producción. \\ \hline

Bagazo &
60 \(^{\circ}\)C, 40 mM &
12.70 &
45 \(^{\circ}\)C, 20 mM &
2.93 &
4.33\(\times\) &
Explotar estabilidad y tolerancia a glucosa. \\ \hline

Biomasa &
70 \(^{\circ}\)C, 80 mM &
36.04 &
50 \(^{\circ}\)C, 40 mM &
7.38 &
4.88\(\times\) &
Operar ENG a mayor temperatura y carga. \\ \hline

Glucósidos &
60 \(^{\circ}\)C, 40 mM &
9.49 &
45 \(^{\circ}\)C, 20 mM &
1.59 &
5.95\(\times\) &
Priorizar producto útil y selectividad. \\ \hline

\end{tabularx}
\end{table}

\subsubsection{A. Alimentos}

La recomendación es operar ENG a \(60\,^{\circ}\mathrm{C}\), 40 mM y ocho horas.
Dentro del rango evaluado produce \(2.78\times\) el mejor resultado de WT, es
decir, aproximadamente 178\% más producto.

La razón es la duración: ENG mantiene actividad durante más tiempo. No conviene
elegir automáticamente \(65\,^{\circ}\mathrm{C}\) por ser la temperatura de
máxima actividad instantánea; el mejor resultado acumulado aparece a
\(60\,^{\circ}\mathrm{C}\).

\subsubsection{B. Bagazo de Caña}

La recomendación es ENG a \(60\,^{\circ}\mathrm{C}\), 40 mM y ocho horas.
La ventaja frente al mejor punto de WT es \(4.33\times\), aproximadamente
333\% más producto.

La prioridad no es aumentar otra vez la velocidad inicial. El valor proviene de
mantener la enzima activa y reducir el efecto de la glucosa acumulada.

\subsubsection{C. Biomasa}

La mejor condición encontrada para ENG es aproximadamente
\(70\,^{\circ}\mathrm{C}\) y 80 mM. El producto alcanza 36.04 unidades relativas,
frente a 7.38 del mejor punto de WT: \(4.88\times\), aproximadamente 388\% más.

La temperatura de máxima actividad instantánea de ENG es
\(75\,^{\circ}\mathrm{C}\), pero el mayor producto acumulado aparece a
\(70\,^{\circ}\mathrm{C}\). Para operación industrial debe priorizarse el segundo
criterio.

\subsubsection{D. Glucósidos}

La recomendación es ENG a \(60\,^{\circ}\mathrm{C}\), 40 mM y ocho horas.
La ventaja frente al mejor punto de WT es \(5.95\times\), aproximadamente
495\% más producto útil.

No conviene aumentar sustrato indefinidamente. La decisión debe equilibrar
inhibición y selectividad, utilizando como objetivo final la cantidad de
glucósido deseado.

\subsection{Prioridad de Explotación}
\label{subsec:prioridad_inversion}

Si el criterio es únicamente el aumento relativo de producto dentro del rango
evaluado, la prioridad es:

\begin{enumerate}
    \item síntesis de glucósidos: \(5.95\times\), aproximadamente \(+495\%\);
    \item procesamiento de biomasa: \(4.88\times\), aproximadamente \(+388\%\);
    \item bagazo de caña: \(4.33\times\), aproximadamente \(+333\%\);
    \item procesamiento de alimentos: \(2.78\times\), aproximadamente \(+178\%\).
\end{enumerate}

Este orden no representa rentabilidad neta. Representa potencial productivo. Para
convertirlo en margen económico real faltan solo tres datos específicos de cada
planta: precio del producto útil, costo de ENG frente a WT y costo energético de
la condición de operación.

La decisión técnica, sin embargo, ya es concreta: glucósidos ofrece la mayor
ventaja relativa; biomasa exige explotar la ventana térmica; bagazo exige
estabilidad y tolerancia; alimentos exige vida operativa.
